\section{Iterative Improvement}\label{appendix:iterative_improvement}
One potential application of \Cref{thm:optimal_stopping} is that when we have a computationally expensive method of evaluating the true reward $\TrueR$, we can design practical training regimes that provably avoid Goodharting. 
Typically training regimes for such reward functions involve iteratively training on a low-cost proxy reward function and fine-tuning on true reward using human feedback~\citep{paulus2017deep}. We can use true reward function to approximate $\Angle{\TrueR}{\ProxR}$, and then optimal stopping gives the optimal amount of time before a ``branching point'' where possible reward training curves diverge (thus, creating Goodharting). 

Specifically, let us assume that we have access to an oracle $\textrm{ORACLE}_{\R^\star}(\R_i, \theta_i)$ which produces increasingly accurate approximations of some true reward $\R^\star$: when called with a proxy reward $\R_i$ and a bound $\theta_i > \Angle{\R^\star}{\R_{i}}$, it returns $\R_{i+1}$ and $\theta_{i+1}$ such that $\theta_{i+1} > \Angle{\R^\star}{\R_{i + 1}}$ and $\lim_{i\to \infty}\theta_i = 0$. Then \cref{alg:iterative-improvement} is an iterative feedback algorithm that avoids Goodharting and terminates at the optimal policy for $\R^\star$. 

\begin{algorithm}
  \caption{Iterative improvement algorithm}\label{alg:iterative-improvement}
  \begin{algorithmic}[1]
   \Procedure{IterativeImprovement}{$\S, \A, \T$}
      \State $\VecR \sim Unif[\RR^{\SxA}]$
      \State $\policy \gets Unif[\RR^{\SxA}]$
      \State $\vec{\occmeaselem}_{-1} = \vec{\occmeaselem}_0 \gets \occmeas{\policy}$
      \State $\theta \gets \frac{\pi}{2}$
      \State $\TangentVec_0 \gets \amax_{\TangentVec \in \ConeTangent(\vec{\occmeaselem}_0)} \TangentVec \cdot \VecR$
      \While{$ \vec{t}_i \neq \vec{0}$}

        \While{$  \vec{\occmeaselem_i} \cdot \TangentVec_i \leq \AngBound$}
          \State $\R, \CosBound \gets$ ORACLE$_{\R^\star}(\R, \CosBound$)
          \State $\VecR \gets \R$
        \EndWhile

        \State $\lambda \gets \max \{\lambda:\vec{\occmeaselem}_i+\lambda \TangentVec_i \in \Polytope\}$
        \State $\vec{\occmeaselem}_{i+1} \gets \vec{\occmeaselem}_i + \lambda \TangentVec_i$
        \State $\TangentVec_{i+1} \gets \amax_{\TangentVec \in \ConeTangent(\vec{\occmeaselem}_{i+1})} \TangentVec \cdot \VecR$
        \State $i \gets i + 1$
      \EndWhile
      \State \Return $\occmeas{\vec{\occmeaselem}_i}^{-1}$
   \EndProcedure
  \end{algorithmic}
\end{algorithm}


\begin{proposition}
\Cref{alg:iterative-improvement} is a valid optimisation procedure, that is, it terminates at the policy $\policy^\star$ which is optimal for the true reward $\R^\star$. 
\end{proposition}
\begin{proof}
By~\Cref{thm:optimal_stopping}, the inner loop of the algorithm maintains that $\J_{\TrueR}(\policy_{i+1}) \geq \J_{\TrueR}(\policy_i)$. If the algorithm terminates, then it must be that $\vec{t_i} = 0$, and the only point that this can happen is in a point $\occmeaselem^{\pi^\star}$. 

Since steepest ascent terminates, showing \cref{alg:iterative-improvement} terminates reduces to showing we only make finitely many calls to $\text{ORACLE}_{\R^\star}$. It can be shown that for any $\TangentVec_i$ produced by steepest ascent on $\R$, $\TangentVec_i = \proj_{\plane}(\R)$ for some linear subspace $\plane$ on the boundary of $\Polytope$ formed by a subset of boundary conditions. Since there are finitely many such $\plane$, there is some $\epsilon$ so that for all $\R, \R^\star$ with $\Angle{\R}{\R^\star} < \epsilon$, $\Angle{\proj_{\plane}(\R)}{\proj_{\plane}(\R^\star)} < \frac{\pi}{2}$ for all $\plane$.

Because we have assumed that $\lim_{i \to \infty} \theta_i = 0$, $\lim_{i \to \infty} \Angle{\R_i}{\R^\star} = 0$ and $\text{ORACLE}_{\R^\star}$ will only be called until $\Angle{\R_i}{\R^\star} = 0$. Then the number of calls is finite and so the algorithm terminates. \end{proof}

We expect this algorithm forms theoretical basis for a  training regime that avoids Goodharting and can be used in practice.
